\section{Results}
\label{sec:results}

\subsection{The model against its own analytic baseline}
\label{sec:headline}

\Cref{tab:headline} gives absolute scores for every configuration on the 429
held-out frames, and \cref{tab:paired} the paired per-frame differences for the
shipped model.

\begin{table}[htbp]
\centering
\caption{Absolute scores on 429 held-out frames. The shipped model (v5) is the
best configuration on degraded input and the worst on clean; only the gated
configuration of \cref{sec:gate} is at or above the analytic baseline on both.}
\label{tab:headline}
\begin{tabular}{llrr}
\toprule
Condition & Method & PU21-PSNR (dB) & CVVDP (JOD) \\
\midrule
clean & analytic baseline            & \textbf{45.99} & 9.448 \\
clean & v5 (32 ch, step \num{81000}) & 42.99 & 9.402 \\
clean & v6 (64 ch)                   & 43.24 & 9.452 \\
clean & \textbf{v5 + shadow gate}    & \textbf{46.06} & \textbf{9.561} \\
\midrule
hard  & analytic baseline            & 25.92 & 7.362 \\
hard  & \textbf{v5}                  & \textbf{27.34} & \textbf{7.805} \\
hard  & v6                           & 26.88 & 7.706 \\
hard  & v5 + shadow gate             & 27.15 & 7.751 \\
\bottomrule
\end{tabular}
\end{table}

\begin{table}[htbp]
\centering
\caption{Paired per-frame difference, v5 against the analytic baseline, 429
frames. The clean PU21 loss of \SI{3.0}{\decibel} corresponds to
$-0.046$~JOD, far below a just-noticeable difference.}
\label{tab:paired}
\begin{tabular}{lrrrrr}
\toprule
 & mean $\Delta$ PU21 & median $\Delta$ & frames won & mean $\Delta$ JOD & median $\Delta$ JOD \\
\midrule
hard  & \textbf{+1.427 dB} & +0.609 & \textbf{348 / 429 (81\%)} & \textbf{+0.443} & +0.134 \\
clean & $-2.999$ dB        & $-2.131$ & 115 / 429 (27\%)        & \textbf{$-0.046$} & $-0.040$ \\
\bottomrule
\end{tabular}
\end{table}

\subsection{Fidelity penalty, perceptual indifference}
\label{sec:disagreement}

\begin{figure}[htbp]
\centering
\includegraphics[width=\linewidth]{fig3_metric_disagreement.pdf}
\caption{Per-frame gain over the analytic baseline, 429 held-out frames;
vertical rules mark the means. The clean distribution is broad and left-shifted
in PU21-PSNR and collapses onto zero in CVVDP. The same frames, the same models,
two metrics, one calling the difference substantial and the other calling it
invisible.}
\label{fig:disagreement}
\end{figure}

A \SI{3.0}{\decibel} PU21-PSNR loss corresponds to $-0.046$~JOD
(\cref{fig:disagreement}), far below a just-noticeable difference. What the
model adds to well-graded input is highlight energy that PU21-PSNR punishes and
that no viewer sees. The clean PU21 row should never be reported without the JOD
beside it. \Cref{sec:expandnet} shows the converse case, in which the two agree
decisively, and \cref{sec:selection} shows the disagreement reappearing inside a
selection rule that touches neither metric.

\subsection{A published method on the same split}
\label{sec:expandnet}

Every number above is against our own analytic baseline, which has been this
paper's largest gap since the first draft. It is now closed for one method.
ExpandNet \citep{marnerides2018} was run on the same 429 frames using the
authors' released weights rather than a reimplementation, with their
pre-processing and post-processing unchanged, on the exact 8-bit frames our model
was given (\cref{tab:expandnet}).

\begin{table}[htbp]
\centering
\caption{Clean condition, 429 frames. ExpandNet is \SI{-18.56}{\decibel} and
$-2.029$~JOD against our deployed configuration, winning 1 of 429 frames on
PU21 and 16 on CVVDP.}
\label{tab:expandnet}
\begin{tabular}{lrr}
\toprule
Method (clean, 429 frames) & PU21 (dB) & CVVDP (JOD) \\
\midrule
\textbf{RUDRA + gate, as deployed} & \textbf{46.06} & \textbf{9.561} \\
RUDRA + gate, exposure-aligned     & 46.61 & 9.558 \\
analytic inverse-ACES baseline     & 45.99 & 9.448 \\
v5 backbone, no gate               & 42.99 & 9.402 \\
ExpandNet                          & 27.50 & 7.532 \\
\bottomrule
\end{tabular}
\end{table}

\paragraph{The exposure alignment, and what it is worth to each side.}
ExpandNet ends in a sigmoid and its released post-processing min-max normalises
each frame to $[0,1]$, so it predicts \emph{relative} radiance with no absolute
anchor. Scoring that directly against a reference in \si{\nits} would measure its
exposure guess. We therefore fit one global scalar per frame on the pixels the
SDR input neither clipped nor crushed. That is a free parameter and it is not
free to grant: ExpandNet's fitted scale has a \textbf{$35.5\times$} spread across
frames (p10 \num{0.49}, p90 \num{17.45}), where ours is $1.0\times$ (p10
\num{0.9993}, p90 \num{1.018}) because our model predicts absolute
\si{\nits}. Put through the same fit our model gains \SI{0.55}{\decibel} and
\emph{loses} \num{0.0025}~JOD, so the row we report for ourselves is the
unaligned one.

The alignment is not what costs ExpandNet its score. Given the \emph{best
possible} per-frame exposure, chosen to maximise its own PU21, it recovers
\SI{+0.39}{\decibel}.

\paragraph{Three caveats, because a gap this size invites the question.}
ExpandNet is run outside its training domain: our SDR frames come from an ACES
approximation at a \SI{-1}{EV} offset, not the curve its corpus used. Our
reference is unclamped and reaches roughly a million \si{\nits}, a range no
method trained on display-referred targets was built for. And ExpandNet
compresses dynamic range by a median factor of $2.3\times$ against the reference
where ours is within $1.32\times$, which contributes to but does not explain
\SI{18}{\decibel}. The row should be read as \emph{this method, unmodified, on
this corpus}, not as a ranking of the field. Santos et
al.~\citep{santos2020}, Eilertsen et al.~\citep{eilertsen2017} and Liu et
al.~\citep{liu2020} have public code and have not been run.

\paragraph{What this settles about \cref{sec:disagreement}.}
The two metrics part company on our model against its own baseline and agree
decisively here: \SI{18.6}{\decibel} and \num{2.03}~JOD say the same thing, and
\num{2}~JOD is two just-noticeable differences. The divergence is a property of
\emph{small} differences on well-graded input, not a defect in either
instrument. When the difference is real, both see it.

\subsection{Capacity is not the constraint}
\label{sec:capacity}

v6 quadruples width to \num{4770117} parameters. It is better on clean
(\SI{+0.25}{\decibel}, $+0.05$~JOD) and worse on hard (\SI{-0.46}{\decibel},
$-0.10$~JOD) than v5: movement in both directions, and smaller than the spread
between conditions. \Cref{sec:levers} places this beside the corpus and
objective ablations.
